\documentclass[10pt]{article}
\usepackage[T1]{fontenc}
\usepackage[utf8]{inputenc}
\usepackage[french]{babel}
\usepackage{amsmath,amssymb,mathrsfs,mathpazo}
\usepackage{geometry}
\geometry{a4paper,margin=2cm}
\usepackage{graphicx,array,tabularx,tabularray,colortbl,multirow,multicol,booktabs,enumitem,pifont,ulem,fancybox,diagbox,cancel}
\usepackage{pstricks,pstricks-add,pst-plot,pst-tree,pst-node,pst-text,pst-func,pst-math,pst-bspline,pst-3dplot,pst-eucl,pst-solides3d}
\let\pscalculate\undefined
\usepackage{pst-bezier}
\usepackage{tikz,pgfplots,tkz-tab}
\usetikzlibrary{arrows,arrows.meta,patterns,calc,positioning,trees,decorations.markings,shapes.geometric,shapes.misc,shapes.arrows,intersections,decorations.pathmorphing,decorations.pathreplacing,backgrounds,shadows,babel}
\pgfplotsset{compat=1.18}
\usepgfplotslibrary{fillbetween}
\usepackage[np]{numprint}
\usepackage{listings,makecell,pas-tableur,twoopt}
\let\barre\undefined
\usepackage[pstricks]{tablvar}
\definecolor{darkgreen}{RGB}{0,100,0}
\providecommand{\red}{\color{red}}
\providecommand{\blue}{\color{blue}}
\providecommand{\green}{\color{green}}
\providecommand{\black}{\color{black}}
\providecommand{\white}{\color{white}}
\providecommand{\gray}{\color{gray}}
\providecommand{\cyan}{\color{cyan}}
\providecommand{\magenta}{\color{magenta}}
\providecommand{\yellow}{\color{yellow}}
\providecommand{\violet}{\color{violet}}
\providecommand{\orange}{\color{orange}}
\providecommand{\euro}{\texteuro}
\providecommand{\EUR}{\texteuro}
\providecommand{\ieme}{\textsuperscript{e}}
\providecommand{\degres}{\ensuremath{^\circ}}
\providecommand{\degree}{\ensuremath{^\circ}}
\providecommand{\vv}[1]{\overrightarrow{#1}}
\providecommand{\Coord}[1]{#1}
\providecommand{\decosix}{}
\providecommand{\starredbullet}{\ensuremath{\bullet}}
\pagestyle{empty}
\setlength{\parindent}{0pt}
\usepackage{adjustbox}\DeclareUnicodeCharacter{2236}{:}\DeclareUnicodeCharacter{2190}{\ensuremath{\leftarrow}}
\providecommand{\R}{\mathbb{R}}
\providecommand{\N}{\mathbb{N}}
\providecommand{\Z}{\mathbb{Z}}
\providecommand{\Q}{\mathbb{Q}}
\providecommand{\D}{\mathbb{D}}
\providecommand{\C}{\mathbb{C}}
\providecommand{\e}{{\mathrm{e}}}
\providecommand{\vect}[1]{\overrightarrow{#1}}
\providecommand{\vectt}[1]{\overrightarrow{\mathrm{#1}}}
\providecommand{\vvt}[1]{\overrightarrow{\mathrm{#1}}}
\providecommand{\np}[2][]{{#2~#1}}
\providecommand{\vv}[1]{{\overrightarrow{#1}}}
\renewcommand{\d}{{\mathrm{d}}}
\providecommand{\degres}{{^\circ}}
\providecommand{\degre}{{^\circ}}
\providecommand{\blue}{\color{blue}}
\providecommand{\red}{\color{red}}
\providecommand{\enskip}{\quad}
\providecommand{\llbracket}{\lbrack\!\lbrack}
\providecommand{\rrbracket}{\rbrack\!\rbrack}
\providecommand{\euro}{{\text{€}}}
\providecommand{\no}{{\text{n°}}}
\providecommand{\interval}[2]{{[#1;#2]}}
\providecommand{\intervFO}[2]{{[#1;#2[}}
\providecommand{\dotfill}{\ldots}

\providecommand{\ds}{\displaystyle}
\providecommand{\textding}[1]{#1}
\providecommand{\Oij}{\(\mathrm{O};\vec{\imath},\vec{\jmath}\)}
\providecommand{\Oijk}{\(\mathrm{O};\vec{\imath},\vec{\jmath},\vec{k}\)}
\providecommand{\Ouv}{\(\mathrm{O};\vec{u},\vec{v}\)}
\providecommand{\intervFF}[2]{\ensuremath{[#1;#2]}}
\renewcommand{\euro}{\ifmmode\text{\texteuro}\else\texteuro\fi}
\DeclareUnicodeCharacter{00B0}{\ifmmode^{\circ}\else\textdegree\fi}
\DeclareUnicodeCharacter{20AC}{\ifmmode\text{\texteuro}\else\texteuro\fi}
% Énoncés : Éducation nationale. Transcription et corrigés : APMEP, Denis Vergès et contributeurs. Découpage et classement : CoopMaths, Sébastien Lozano, Éric Elter et contributeurs. Adaptation de la présentation : Maths Explorer.
% CC BY-SA : https://coopmaths.fr/www/about/
% https://coopmaths.fr/alea/static/eam/2026/tex/eam_specifique_2026_06_asie_3.tex
\begin{document}
\noindent\textit{APMEP / CoopMaths. CC BY-SA. Présentation adaptée par Maths Explorer.}\par\medskip
Cet exercice s'intéresse à la croissance d'une personne appelée Juliette. Sa croissance est modélisée par la fonction $f$, représentée ci-dessous.

La \textbf{vitesse de croissance} de Juliette (en cm par an) à une date donnée $t$ est égale au \textbf{nombre dérivé} $\boldsymbol{f'(t)}$.

On a également représenté trois tangentes à la courbe représentative de $f$ : celles aux points A, B et C.

\begin{center}
\begin{tikzpicture}[x = 1cm, y=1mm]
	\draw [color = gray!75, line width = 0.5pt, xstep=1, ystep = 10] (5,90) grid (19,180);
	\draw[<->, line width=1pt,>=stealth] (19,90)--(5,90)node[pos=0.5, below=1cm]{Âge (en années)} -- (5,180) node[pos=0.5,sloped,above=1cm]{Taille (en cm)};
	\foreach \x in {5,...,19}
	\draw[shift={(\x,90)},color=black] (0pt,2pt) -- (0pt,-2pt) node[below, fill = white] {\footnotesize $\np{\x}$};
	\foreach \y in {90,100,...,180}
	\draw[shift={(5,\y)},color=black] (2pt,0pt) -- (-2pt,0pt) node[left, fill = white] {\footnotesize $\np{\y}$};
	\clip (5,90) rectangle (19,180);
	\draw[dashed, line width=1pt,teal] (5,110)--(17,182)
		(5,137.5)--(19,172.5)
		(5,163)--(19,163) ;
	\draw[line width=1.2pt,color=red,smooth]
	(5,109) ..controls (7,118) and (9,132)..
	(10,139)..controls (11,145) and (11,146)..
	(12,152)..controls (12.5,155) and (13.5,158.75)..
	(14,160)..controls (14.5,161.25) and (16.5,163)..
	(17,163) -- (19,163);
	\node[red] at (7.5,117) {$\mathcal{C}_f$};
	\foreach \nom/\coord/\pos in {A/{(12,152)}/{below right},B/{(14,160)}/{below right},C/{(17,163)}/{above right}}
		{\draw[shift={\coord}] (-2pt,-2pt)--(2pt,2pt) (-2pt,2pt)--(2pt,-2pt) (0,0) node[\pos]{\nom} ;}
\end{tikzpicture}
\end{center}

On répondra aux questions suivantes avec la précision permise par le graphique :

	\begin{enumerate}
		\item Estimer la taille de Juliette le jour de ses 14 ans.

		\item À quel âge Juliette grandissait-elle le plus vite, à 12 ans ou à 14 ans ? \emph{Justifier.}

		\item Pourquoi peut-on considérer que Juliette a terminé sa croissance à 17 ans ?

		\item Estimer la taille adulte de Juliette.
	\end{enumerate}

\end{document}
